\scanpage{23}
Il s'impose de regarder \unsure{égalmt|égalt}, si $A$ est ordinaire
\[
T_p(A)_{\mathrm{\acute et}}=\varprojlim\deleted{\unsure{T_p}}({}_{p^{\unsure{\nu|n}}}A)_{\mathrm{\acute et}}
\]
(N.B. $T_p(A)$ est une extension
\[
0\to T_p(A)_{\mathrm{rad}}\to T_p(A)\mathbin{\to}T_p(A)_{\mathrm{\acute et}}\to0
\]
On sait, par dualité de Cartier, que dans la
suite exacte correspondante
\[
0\mathbin{-}T_p(A^*)_{\mathrm{rad}}\mathbin{-}T_p(A^*)\mathbin{-}T_p(A^*)_{\mathrm{\acute et}}\mathbin{-}0
\]
les termes médians sont duals l'un de
l'autre, et \deleted{plus précisément} cette dualité est
compatible avec \unsure{filtration|filtrations}, donc
\[
\begin{aligned}
T_p(A)_{\mathrm{rad}}&\simeq D\bigl(T_p(A^*)_{\mathrm{\acute et}}\bigr)\\
T_p(A^*)_{\mathrm{rad}}&\simeq D\bigl(T_p(A)_{\mathrm{\acute et}}\bigr)
\end{aligned}
\]
donc la connaissance de $T_p(A)_{\mathrm{\acute et}}$ et $T_p(A^*)_{\mathrm{\acute et}}$
permet de reconstituer les \uline{deux} facteurs
de $T_p(A)$, donc $T_p(A)$ si la base est
parfaite [en général, il y aurait lieu d'
expliciter la structure d'extension \ldots].

On aurait envie d'établir un isom can.
\deleted{fonctoriel}
% archive sheet number: 22
